\documentclass[12pt]{article}
\usepackage[spanish]{babel}
\usepackage[utf8]{inputenc}
\usepackage[T1]{fontenc}
\usepackage{geometry}
\usepackage{graphicx}
\usepackage{hyperref}
\geometry{margin=2.5cm}

\title{Laboratorio Practico N2: Implementacion de Stack Web con Docker}
\author{Tomas Teihuel}
\date{\today}

\begin{document}
\maketitle

\section{Introduccion}
Este informe presenta el despliegue de una aplicacion web dinamica con operaciones CRUD usando Docker. La solucion utiliza un contenedor PHP con Apache para la aplicacion y un contenedor MariaDB para la base de datos.

\section{Desarrollo}
La arquitectura se compone de dos servicios definidos en \texttt{docker-compose.yml}. El servicio \texttt{app} se construye desde un \texttt{Dockerfile} propio y se conecta al servicio \texttt{db} mediante una red virtual de Docker. La base de datos usa un volumen persistente llamado \texttt{db\_data}.

\subsection{Construccion}
Agregar captura del comando:
\begin{verbatim}
docker compose build
\end{verbatim}

\subsection{Ejecucion}
Agregar captura de los comandos:
\begin{verbatim}
docker compose up -d
docker ps
\end{verbatim}

\subsection{Logs y red}
Agregar captura del comando:
\begin{verbatim}
docker logs laboratorio_db
\end{verbatim}

\subsection{Prueba funcional}
Agregar captura de la aplicacion funcionando en:
\begin{verbatim}
http://localhost:8080
\end{verbatim}

\section{Configuracion Docker}
El contenedor de aplicacion usa la imagen \texttt{php:8.3-apache} y habilita \texttt{pdo\_mysql}. El contenedor de base de datos usa la imagen oficial \texttt{mariadb:11.4}. Las credenciales se cargan desde variables de entorno y la informacion de la base de datos persiste mediante un volumen.

\section{Conclusion}
La implementacion permite desplegar una aplicacion CRUD de forma reproducible usando Docker Compose. La separacion entre aplicacion y base de datos facilita la administracion, el despliegue y la explicacion tecnica de redes, volumenes e imagenes.

\section{Anexos}
\begin{itemize}
    \item Link repositorio GitHub: \url{https://github.com/TomasTeihuel/lab-docker-crud}
    \item Link documento Overleaf: agregar URL.
    \item Archivos principales: \texttt{Dockerfile}, \texttt{docker-compose.yml}, \texttt{init.sql}, \texttt{README.md}.
\end{itemize}

\end{document}
